\section{Estructura e Implementación de la Red Neuronal Generativa}

El núcleo generativo del sistema de composición automatizada se fundamenta en una red neuronal artificial híbrida construida y compilada de manera paramétrica mediante la función \texttt{build\_dynamic\_model} en el módulo \texttt{constructor\_red.py}. La entrada del modelo está diseñada para consumir tensores secuenciales de dimensiones $(L, F)$, donde $L$ (SEQ\_LENGTH) representa la longitud de la ventana temporal (fijada en 4 notas o eventos musicales previos) y $F$ (NUM\_FEATURES) representa la dimensión del vector de características por nota (fijada en 4 parámetros: clase de altura, duración rítmica, grado de la escala diatónica activa e intervalo en tonos respecto a la nota previa).

La arquitectura profunda está organizada en una topología secuencial secuenciada a través de las siguientes capas especializadas:
\begin{enumerate}
    \item \textbf{Capa de Convolución Temporal (Conv1D):} Recibe el tensor de entrada y actúa como un extractor de características locales a nivel de sub-frase musical. Utiliza 32 filtros y un tamaño de kernel de 2. Esta capa extrae patrones intervalares y de contorno rítmico inmediatos antes de alimentar las unidades recurrentes.
    \item \textbf{Capa de Memoria a Corto y Largo Plazo (LSTM):} Compuesta por 64 unidades con función de activación ReLU. Esta capa recurrente captura las dependencias temporales y la memoria formal a largo plazo de la secuencia musical, permitiendo mantener la coherencia armónica a lo largo de los compases. Si se definen múltiples capas recurrentes en la matriz, el sistema infiere y activa de forma automática el parámetro \texttt{return\_sequences=True} en las capas previas.
    \item \textbf{Capa de Regularización (Dropout):} Aplica una tasa de descarte aleatorio de conexiones de $0.20$ ($20\%$). Esta capa actúa como un regulador indispensable para evitar el sobreajuste (overfitting) del modelo ante la memorización directa del conjunto de entrenamiento.
    \item \textbf{Capa Densa Intermedia (Dense):} Una capa totalmente conectada compuesta por 32 neuronas con activación ReLU, la cual actúa como un extractor de representaciones abstractas finales antes de la síntesis de salida.
    \item \textbf{Capa de Salida (Dense con Regresión):} A diferencia de los modelos clasificadores de lenguaje tradicionales, esta capa final consta de un número de unidades igual a \texttt{NUM\_FEATURES} (4 neuronas) con una función de activación lineal. Esto convierte el problema en una regresión lineal multivariada, estimando valores continuos para los cuatro atributos de la siguiente nota.
\end{enumerate}

El modelo se compila utilizando el optimizador Adam, el cual aplica tasas de aprendizaje adaptativas para acelerar la convergencia, y la función de pérdida del error cuadrático medio (MSE), alineada con la naturaleza continua del espacio métrico y de alturas. Este diseño permite que la red proponga melodías fluidas e inéditas basadas en la herencia estilística de los datos de entrada, las cuales son posteriormente normalizadas y corregidas por la capa del árbol de decisión determinista.
